\chapter{Concluzii și perspective de dezvoltare}

\section{Concluzii}
Prin prezenta lucrare am abordat dezvoltarea unei aplicații de planificare a călătoriilor, concepută pentru a facilita procesul organizării unui itinerariu complex. Aplicația integrează cu succes mai multe tehnologii și caracteristici într-o platformă care ușurează experiența de planificare.

\begin{itemize}
    \item Dezvoltarea unui algoritm care colectează preferințele utilizatorilor și generează itinerarii personalizate pe baza acestora. Această abordare depășește simplele sugestii de „cele mai vizitate atracții” și oferă o experiență personalizată fiecărui utilizator.
    \item Centralizarea tuturor informațiilor (rezervări, cazări, locații) într-o singură platformă. Acest lucru elimină dificultatea jonglării cu mai multe aplicații și a gestionării logisticii.
    \item Capacitatea sistemului de a actualiza itinerariile, permițând utilizatorilor să bifeze locațiile vizitate și să regenereze planurile rămase pe baza progresului real. Replanificarea asigură o planificare optimă a traseului pe tot parcursul călătoriei.
\end{itemize}

\section{Perspective de dezvoltare}
În ceea ce privește perspectivele de dezvoltare, aplicația are numeroase posibilități de extindere:

\begin{itemize}
    \item Integrarea modelelor de învățare automată pentru a prezice timpii optimi de călătorie pe baza datelor istorice sau pe baza traficului în timp real, oferind astfel o planificare și mai precisă.
    \item Analiza predictivă pentru a sugera destinații pe baza modelelor de comportament ale utilizatorilor și a experiențelor similare de călătorie.
    \item Planificarea colaborativă pentru scenarii de călătorie în grup, cu recomandări comune.
    \item Integrarea cu servicii meteo pentru prognoza meteo în timp real și recomandări de călătorie.
    \item Servicii bazate pe locație și geofencing pentru a oferi memento-uri, informații și autocompletarea itinerariului în timp real.
    \item Integrarea calculării rutelor folosind transportul public sau alte opțiuni de transport pentru a oferi o planificare completă a călătoriei.
\end{itemize}